\originalsection{作业 9: 碰撞}
\begin{exercise}
原题 1. 乒乓球以 $v_{p0}$ 沿 $+x$ 与沿 $-x$ 速率 $v_{b0}$ 的保龄球弹性碰撞, 保龄球质量远大于乒乓球. 求乒乓球末速度的 $x$ 分量.
\end{exercise}
\begin{solution}
保龄球速度近似不变为 $-v_{b0}$. 在其系中乒乓球速度 $v_{p0}+v_{b0}$ 反号, 回地面系得 $v_{pf,x}=-v_{p0}-2v_{b0}$. 初末速率与带符号速度必须区分.
\end{solution}
\begin{exercise}
原题 2. 地球以 $v_e\uv x$ 移动. 质量可忽略的卫星从 $-x$ 方向以 $4v_e\uv x$ 接近, 受地球引力作双曲线飞掠后朝原来来处返回. 求远离后速率.
\end{exercise}
\begin{solution}
地球系入射速度 $3v_e\uv x$, 题设反向出射为 $-3v_e\uv x$, 回原系得 $-2v_e\uv x$, 所以速率 $2v_e$. 能量在地球系守恒, 卫星在原系可与地球交换能量.
\end{solution}
\begin{exercise}
原题 3. 质量 $m_b$ 的块带半径 $R$ 的曲面, 静止于光滑桌面. 小立方体质量 $m_c<m_b$ 以水平速率 $v_{c0}$ 滑入, 全部接触光滑, 大块可运动. (a) 求小体最大升高 $h<R$. (b) 下滑后分离, 求两物体速度大小.
\end{exercise}
\begin{solution}
(a) 到最高相对位置两者水平速度相同, 由动量得 $V=m_cv_{c0}/(m_b+m_c)$. 能量给 $m_cgh=m_cv_{c0}^2/2-(m_b+m_c)V^2/2$, 所以 $h=m_bv_{c0}^2/[2g(m_b+m_c)]$. (b) 回到同高度后是等效一维弹性过程, 小体速度 $v_c=(m_c-m_b)v_{c0}/(m_b+m_c)<0$, 大块 $v_b=2m_cv_{c0}/(m_b+m_c)>0$, 速度大小取绝对值.
\end{solution}
\begin{exercise}
原题 4. 块 $m_1$ 以速率 $v_1$ 沿光滑桌面撞静止的弹簧--块 $m_2$, 弹簧劲度 $k$. 一维碰撞, 求最大压缩量.
\end{exercise}
\begin{solution}
压缩最大时相对速度为零. 质心平动能保留, 初始相对动能 $\mu v_1^2/2$ 全变为弹簧能, $\mu=m_1m_2/(m_1+m_2)$. 所以 $X=v_1\sqrt{\mu/k}$.
\end{solution}
\begin{exercise}
原题 5. 块 $M$ 由弹簧 $k$ 连墙, 在光滑桌面静止且弹簧未形变. 子弹 $m_b$ 以 $v_0$ 射入并粘住. (a) 碰后速度. (b) 首次停住距离. (c) 碰后与碰前动能比. (d) 改成动摩擦 $\mu_k$ 桌面, 块首次返回初位时速率为碰后的一半, 求这一往返总路程.
\end{exercise}
\begin{solution}
(a) $V=m_bv_0/(M+m_b)$. (b) $X=V\sqrt{(M+m_b)/k}$. (c) 比为 $m_b/(M+m_b)$. (d) 起终弹簧能同为零, 摩擦耗散 $\mu_k(M+m_b)gD=\frac12(M+m_b)[V^2-(V/2)^2]$, 得 $D=3V^2/(8\mu_kg)$. 该实测末速已编码弹簧与摩擦的共同作用, 因而结果可不含 $k$.
\end{solution}
\begin{exercise}
原题 6. $m$ 粒子沿 $+x$ 以 $v_i$ 撞静止的 $2m$ 粒子, 完全弹性. 后者以与 $+x$ 夹角 $45^\circ$ 向右上方出射, 前者向右下方出射, 图中 $\theta_1$ 从 $+x$ 顺时针量. 求 $\tan\theta_1$.
\end{exercise}
\begin{solution}
设重粒子末速 $w$, 则轻粒子分量 $v_{1x}=v_i-\sqrt2w$, $v_{1y}=-\sqrt2w$. 能量给 $v_i^2=(v_i-\sqrt2w)^2+2w^2+2w^2$, 非零碰撞解 $w=\sqrt2v_i/3$. 因而 $v_{1x}=v_i/3$, $v_{1y}=-2v_i/3$, 图示正角满足 $\tan\theta_1=2$.
\end{solution}
